% ===========================================================================
%  Chapter 2: Related Work
%  Dual-Target RL Molecular Generation — Workshop Report
% ===========================================================================

This section surveys six areas of prior work that motivate and contextualize our
dual-target reinforcement-learning pipeline: multi-target drug design
(Section~\ref{sec:rw:mtdd}), de~novo generation with RL
(Section~\ref{sec:rw:rlgen}), multi-objective optimization
(Section~\ref{sec:rw:moo}), docking-based scoring
(Section~\ref{sec:rw:docking}), learned affinity prediction
(Section~\ref{sec:rw:boltz2}), and clinical validation of AI-designed
therapeutics (Section~\ref{sec:rw:clinical}).

% ---------------------------------------------------------------------------
\subsection{Multi-Target Drug Design and Polypharmacology}
\label{sec:rw:mtdd}
% ---------------------------------------------------------------------------

Complex diseases rarely arise from single-gene perturbations.
Alzheimer's disease (AD), for instance, involves amyloid precursor processing
(BACE1), tau phosphorylation cascades (GSK3$\beta$, CDK5), and
cholinergic deficit (AChE, BChE), motivating simultaneous modulation of
multiple targets~\citep{bolognesi2013polypharmacology}.
Polypharmacology---the deliberate design of ligands that engage two or more
targets---offers potential advantages over combination therapy: simplified
pharmacokinetics, reduced drug--drug interactions, and a single
ADMET profile to optimize~\citep{stern2022multitarget}.
Clinical precedents include ladostigil, a dual AChE/MAO-B inhibitor
developed for AD~\citep{weinreb2012cholinesterase}, and donepezil--memantine
fixed-dose combinations.

The computational challenge is non-trivial: binding sites of target pairs
such as BACE1/GSK3$\beta$---where GSK3$\beta$ directly regulates BACE1
expression~\citep{ly2013gsk3b}---share little structural homology, so
optimizing affinity to one pocket often degrades binding to the other.
Notably, \citet{jiang2020notopterol} identified Notopterol as a natural
dual BACE1/GSK3$\beta$ inhibitor via traditional screening; our work
instead uses RL-based generation to explore the joint chemical space
systematically.

% ---------------------------------------------------------------------------
\subsection{De Novo Molecular Generation with Reinforcement Learning}
\label{sec:rw:rlgen}
% ---------------------------------------------------------------------------

\paragraph{Fragment-based RL.}
FREED~\citep{yang2021freed} introduced a fragment-attachment MDP where an
agent iteratively selects molecular fragments to extend a growing molecule,
receiving a docking-score reward.
FREED++~\citep{telepov2023freedpp} provided a thorough reproduction and
improvement: explorative experience replay, reward shaping, and stabilized
training yielded consistent gains on standard benchmarks.
Our pipeline adopts FREED++ as the base generator.
In \texttt{double} mode, we monkey-patch the reward computation to call
UniDock on both targets per RL step, combining scores via a worst-case
scalarization ($R = -\max(d_A, d_B)$) before returning the reward signal
to the agent.  This is, to our knowledge, the first dual-target extension
of FREED++.

\paragraph{SMILES-based RL.}
REINVENT~4~\citep{loeffler2024reinvent4} is the current release of the
AstraZeneca generative platform, supporting de~novo, R-group, linker,
and Mol2Mol modes with a modular scoring stack.
It uses a DAP (Difference between Agent and Prior) loss and supports
curriculum learning.
ADMETrix~\citep{mourdoukoutas2025admetrix} extends REINVENT with 27 ADMET
and 10 physicochemical predictors via the ADMET-AI model, demonstrating
that RL agents can optimize pharmacological profiles beyond single
binding scores.
While REINVENT employs weighted-sum scalarization for multi-objective
tasks, our design deliberately avoids this choice to isolate the
dual-target effect under a single, fixed architecture.

\paragraph{Synthesis-aware generation.}
A recent line of work constrains generation to synthetically accessible
molecules.
SynFlowNet~\citep{cretu2025synflownet} applies GFlowNets in a
reaction-template space, jointly producing molecules and synthesis
routes.
SynGFN~\citep{kim2025syngfn} introduces hierarchical pre-trained policies
over cascaded simulated reactions.
RxnFlow~\citep{lee2025rxnflow} scales to 1.2\,M building blocks and 71
reaction templates, achieving state-of-the-art mean Vina scores
($-8.85$\,kcal/mol on CrossDocked2020).
These approaches are complementary to ours: our focus is the
dual-target reward signal rather than synthetic feasibility, though
combining both is a promising future direction.

% ---------------------------------------------------------------------------
\subsection{Multi-Objective Optimization in Molecular Design}
\label{sec:rw:moo}
% ---------------------------------------------------------------------------

Multi-objective molecular optimization seeks to balance competing
properties---binding affinity to multiple targets, synthetic accessibility,
ADMET constraints---that cannot be collapsed into a single scalar without
loss of information.
\citet{fromer2023moo} provide a comprehensive taxonomy of approaches,
distinguishing \emph{scalarization} methods (which reduce the problem to
a single objective) from \emph{Pareto-based} methods (which maintain a
population on the Pareto front).

\paragraph{Scalarization.}
Weighted sums are the simplest and most widely used
approach~\citep{loeffler2024reinvent4, liu2023drugexv3}; however, they
can only recover convex regions of the Pareto front.
Tchebycheff decomposition~\citep{zhang2007moead} overcomes this limitation
and is used in DrugEx~v3~\citep{liu2023drugexv3}.
Our primary reward, $\max(d_A, d_B)$ (worst-case), is a special case of
Tchebycheff with equal weights and a zero reference point: it penalizes
molecules that sacrifice one target for the other, providing a
conservative lower bound on dual-target quality.

\paragraph{Pareto-based methods.}
DrugEx~v2~\citep{liu2021drugex} pioneered Pareto-ranked RL for
polypharmacological design, using non-dominated sorting within a
population-based policy.
ParetoDrug~\citep{paretodrug2024} combines Pareto MCTS with a pre-trained
autoregressive model, using a ParetoPUCT criterion to balance exploration
and exploitation.
PMMG~\citep{pmmg2025} extends Pareto MCTS to seven simultaneous
objectives (dual EGFR/HER2 affinity, QED, SA, and toxicity), achieving
51.65\% success rate---2.5$\times$ higher than prior work.
CMOMO~\citep{cmomo2025} introduces dynamic constraints in a two-stage
framework, doubling the success rate for GSK3$\beta$ inhibitors.
ScaRL-P~\citep{wang2024scarlp} applies clustered Pareto RL with
per-objective critics.

\paragraph{Quality metrics.}
The \emph{hypervolume indicator} measures the volume of objective space
dominated by a Pareto front, providing a single scalar that captures both
convergence and diversity~\citep{verhellen2022hypervolume}.
HMAMP~\citep{wang2025hmamp} uses hypervolume-driven search for
antimicrobial peptide design; we adopt the hypervolume as a secondary
endpoint (Section~5).

\paragraph{Our positioning.}
We restrict the primary experiment to a single scalarization
($\max(d_A, d_B)$) rather than employing Pareto-based methods.  This isolates the effect of \emph{dual-target RL generation}
under controlled conditions: identical generator (FREED++), identical
molecular budget, and identical reward architecture, varying only the
number of targets.  Pareto and hypervolume alternatives are evaluated
offline as sensitivity analyses.

% ---------------------------------------------------------------------------
\subsection{Docking-Based Scoring Functions}
\label{sec:rw:docking}
% ---------------------------------------------------------------------------

AutoDock Vina~\citep{trott2010vina} remains the most widely used
open-source docking program, employing an empirical scoring function
that combines steric, hydrophobic, and hydrogen-bond terms with an
iterated local search optimizer.
The CASF-2016 benchmark~\citep{su2019casf2016} reported Vina's scoring
power as Pearson $r = 0.604$ and ranking power as Spearman
$\rho = 0.528$---competitive among classical scoring functions but far
from perfect.

UniDock~\citep{yu2023unidock} ports the Vina scoring function to GPU,
achieving up to 1000$\times$ speedup over CPU Vina while preserving
identical scoring semantics.
This makes UniDock practical as an RL reward: each molecule can be
docked in milliseconds, enabling the thousands of evaluations required
per RL epoch.
Our pipeline uses UniDock GPU (\texttt{unidock\_gpu}) as the docking
backend for both single- and dual-target reward computation.

\paragraph{Known limitations.}
Docking scores are approximate: they neglect explicit solvation,
protein flexibility, and entropic contributions.
PoseX~\citep{posex2025} demonstrated that on certain targets, Vina-family
scoring functions perform near-randomly for virtual screening hit
identification.
Performance is target-dependent: some pockets are well-suited to
empirical scoring, while others require more sophisticated methods.
These limitations motivate our use of an orthogonal evaluation
surrogate (Boltz-2, Section~\ref{sec:rw:boltz2}) and our per-target
Spearman $\rho$ sanity check between docking scores $d_X$ and
learned affinities $a_X$.

% ---------------------------------------------------------------------------
\subsection{Learned Affinity Prediction: Boltz-2}
\label{sec:rw:boltz2}
% ---------------------------------------------------------------------------

Boltz-2~\citep{passaro2025boltz2} is a unified biomolecular prediction
model that extends a structure-prediction backbone (PairFormer) with
binding-likelihood and affinity-prediction heads.
The affinity module is trained on approximately 750\,K compound--target
pairs drawn from ChEMBL, BindingDB, and PDBbind, spanning Ki, Kd,
IC$_{50}$, and EC$_{50}$ assay types.
The training loss combines Huber regression on absolute values with
pairwise intra-assay differences, encouraging consistent ranking
within each assay context.
The output is $\log_{10}(\text{IC}_{50}\text{-like},\;\mu\text{M})$,
where lower values indicate stronger predicted binding; notably, this
is \emph{not} pIC$_{50}$ (which is $-\log_{10}$).

\paragraph{Performance.}
On the authors' four-target FEP+ validation subset (CDK2, TYK2, JNK1,
p38), Boltz-2 achieves Pearson $r = 0.66$, comparable to OpenFE and
below commercial FEP+ ($r = 0.78$).
Independent evaluations report more modest correlations: the PL-REX
benchmark (Yesylevskyy) yields mean $r \approx 0.42$, while DeepMirror
per-target evaluations range from $r = 0.26$ (cGAS, narrow activity
range) to $r = 0.81$ (Mpro).
These results underscore target-dependent performance and motivate our
per-target Spearman $\rho$ sanity check.

\paragraph{Limitations.}
The affinity module does not explicitly model cofactors, ions, or
multimeric binding partners; predictions for cofactor-containing pockets
(e.g., MAO-B with FAD) are unreliable.
The predicted values span a compressed range (typically
$-2$ to $+3$ log units) relative to experimental data, and a
$\sim$40\% false-positive rate at stringent thresholds has been
reported.
We therefore use Boltz-2 exclusively as an \emph{evaluation surrogate}
for relative comparison between generation strategies, never as a
training reward, and interpret hit-rate differences rather than
absolute predicted IC$_{50}$ values.

\paragraph{Ecosystem.}
Several recent tools build on Boltz-2.
Boltzina~\citep{furui2025boltzina} combines AutoDock Vina docking poses
with the Boltz-2 affinity head, achieving up to 11.8$\times$ speedup
over full Boltz-2 inference.
Most relevant to our work, SynFlowNet~\citep{cretu2025synflownet}
uses Boltz-2 as a scoring function within a GFlowNet RL loop
(off-policy, with replay buffer), providing a direct precedent for
Boltz-2 as an RL scoring surrogate in molecular generation.
Our pipeline differs in that Boltz-2 is used post-hoc for evaluation
rather than as the RL reward, avoiding potential mode collapse
toward Boltz-2-specific artifacts.

% ---------------------------------------------------------------------------
\subsection{Clinical Validation of AI-Generated Therapeutics}
\label{sec:rw:clinical}
% ---------------------------------------------------------------------------

While computational molecular design has a long history, recent clinical
milestones demonstrate that AI-generated molecules can reach the clinic.
Rentosertib~\citep{xu2025rentosertib}, a TNIK inhibitor for idiopathic
pulmonary fibrosis, was designed end-to-end using AI: PandaOmics
identified the target, Chemistry42 generated the molecule, and the
program progressed from target identification to preclinical candidate
in 18 months.
In a Phase~2a randomized trial (NCT05938920), rentosertib improved FVC
by $+98.4$\,mL versus $-20.3$\,mL for placebo, providing the first
clinical evidence that an AI-designed molecule can achieve efficacy
in humans~\citep{xu2025rentosertib}.

In the polypharmacology domain, POLYGON~\citep{ren2024polygon}
generated 32 synthesized dual MEK1/mTOR compounds, with the majority
showing $>$50\% inhibition of both targets at 10\,$\mu$M.
Automated multi-target ligand design via chemical language
models~\citep{zhang2024clm} produced 12 synthesized compounds across
six target pairs.
These experimental validations, while still pre-clinical, demonstrate
that multi-target generative design can yield physically real,
biologically active molecules, motivating further computational
investigation of dual-target RL strategies.
